% Erratum E6 for inclusion in main.tex.
% Suggested use: % Erratum E6 for inclusion in main.tex.
% Suggested use: % Erratum E6 for inclusion in main.tex.
% Suggested use: % Erratum E6 for inclusion in main.tex.
% Suggested use: \input{erratum_E6_prop22_endpoint_nu0}

\subsection*{Erratum E6: the endpoint $\nu=0$ in Proposition 2.2}
\label{s:erratum:E6}

In the proof of Proposition 2.2, specifically in the proof of
inequality (26b), the displayed choice
\[
        t=s^\gamma\sqrt{\frac{\gamma-\nu}{\nu}}
\]
is meaningful only for $0<\nu<\gamma$.  The endpoint $\nu=0$ should be
handled separately.

For $0<\nu<\gamma$, the printed argument is valid.  Indeed, putting
$\theta:=\nu/\gamma$, the displayed choice of $t$ is exactly the value
for which
\[
        s^{-2\nu}
        =
        N_{\theta}^{2}t^{-2\theta}
        \frac{t^2}{t^2+s^{2\gamma}} .
\]
For $\lambda\leq s$ this gives
\[
        s^{-2\nu}
        \leq
        N_{\theta}^{2}t^{-2\theta}
        \frac{t^2}{t^2+\lambda^{2\gamma}},
\]
and hence, using equation (24),
\[
        s^{-2\nu}\norm{E_{[0,s)}x}{}^2
        \leq
        N_{\theta}^{2}t^{-2\theta}|K_t^\gamma(x)|^2
        \leq
        N_{\theta}^{2}\norm{x}{\nu:\gamma}^{2} .
\]
Taking the supremum over $s>0$ yields
\[
        \normiii{x}{\nu}
        \leq
        N_{\nu/\gamma}\norm{x}{\nu:\gamma}
\]
for $0<\nu<\gamma$.

It remains only to record the endpoint $\nu=0$.  Since $N_0=1$, the
claim becomes
\[
        \normiii{x}{0}\leq\norm{x}{0:\gamma}.
\]
In fact equality holds.  First,
\[
        \normiii{x}{0}
        =
        \sup_{s>0}\norm{E_{[0,s)}x}{}
        =
        \norm{x}{}.
\]
The last equality follows because the spectral measure of $T^*T$ is
compactly supported and $E_{[0,s)}$ equals the identity for all sufficiently
large $s$.  On the other hand, by equation (24),
\[
        |K_t^\gamma(x)|^2
        =
        \int_{[0,\infty)}
        \frac{t^2}{t^2+\lambda^{2\gamma}}\,d\mu_x(\lambda)
        \leq
        \int_{[0,\infty)}d\mu_x(\lambda)
        =
        \norm{x}{}^2 .
\]
Thus $\norm{x}{0:\gamma}\leq\norm{x}{}$.  Conversely, as
$t\to\infty$ the integrand in equation (24) increases pointwise to $1$;
by monotone convergence,
\[
        \lim_{t\to\infty}|K_t^\gamma(x)|^2
        =
        \norm{x}{}^2 .
\]
Therefore
\[
        \norm{x}{0:\gamma}
        =
        \sup_{t>0}|K_t^\gamma(x)|
        =
        \norm{x}{}
        =
        \normiii{x}{0}.
\]
This proves inequality (26b) also at $\nu=0$.  Thus the statement of
Proposition 2.2 is unchanged; only the endpoint argument needs
to be supplied separately.


\subsection*{Erratum E6: the endpoint $\nu=0$ in Proposition 2.2}
\label{s:erratum:E6}

In the proof of Proposition 2.2, specifically in the proof of
inequality (26b), the displayed choice
\[
        t=s^\gamma\sqrt{\frac{\gamma-\nu}{\nu}}
\]
is meaningful only for $0<\nu<\gamma$.  The endpoint $\nu=0$ should be
handled separately.

For $0<\nu<\gamma$, the printed argument is valid.  Indeed, putting
$\theta:=\nu/\gamma$, the displayed choice of $t$ is exactly the value
for which
\[
        s^{-2\nu}
        =
        N_{\theta}^{2}t^{-2\theta}
        \frac{t^2}{t^2+s^{2\gamma}} .
\]
For $\lambda\leq s$ this gives
\[
        s^{-2\nu}
        \leq
        N_{\theta}^{2}t^{-2\theta}
        \frac{t^2}{t^2+\lambda^{2\gamma}},
\]
and hence, using equation (24),
\[
        s^{-2\nu}\norm{E_{[0,s)}x}{}^2
        \leq
        N_{\theta}^{2}t^{-2\theta}|K_t^\gamma(x)|^2
        \leq
        N_{\theta}^{2}\norm{x}{\nu:\gamma}^{2} .
\]
Taking the supremum over $s>0$ yields
\[
        \normiii{x}{\nu}
        \leq
        N_{\nu/\gamma}\norm{x}{\nu:\gamma}
\]
for $0<\nu<\gamma$.

It remains only to record the endpoint $\nu=0$.  Since $N_0=1$, the
claim becomes
\[
        \normiii{x}{0}\leq\norm{x}{0:\gamma}.
\]
In fact equality holds.  First,
\[
        \normiii{x}{0}
        =
        \sup_{s>0}\norm{E_{[0,s)}x}{}
        =
        \norm{x}{}.
\]
The last equality follows because the spectral measure of $T^*T$ is
compactly supported and $E_{[0,s)}$ equals the identity for all sufficiently
large $s$.  On the other hand, by equation (24),
\[
        |K_t^\gamma(x)|^2
        =
        \int_{[0,\infty)}
        \frac{t^2}{t^2+\lambda^{2\gamma}}\,d\mu_x(\lambda)
        \leq
        \int_{[0,\infty)}d\mu_x(\lambda)
        =
        \norm{x}{}^2 .
\]
Thus $\norm{x}{0:\gamma}\leq\norm{x}{}$.  Conversely, as
$t\to\infty$ the integrand in equation (24) increases pointwise to $1$;
by monotone convergence,
\[
        \lim_{t\to\infty}|K_t^\gamma(x)|^2
        =
        \norm{x}{}^2 .
\]
Therefore
\[
        \norm{x}{0:\gamma}
        =
        \sup_{t>0}|K_t^\gamma(x)|
        =
        \norm{x}{}
        =
        \normiii{x}{0}.
\]
This proves inequality (26b) also at $\nu=0$.  Thus the statement of
Proposition 2.2 is unchanged; only the endpoint argument needs
to be supplied separately.


\subsection*{Erratum E6: the endpoint $\nu=0$ in Proposition 2.2}
\label{s:erratum:E6}

In the proof of Proposition 2.2, specifically in the proof of
inequality (26b), the displayed choice
\[
        t=s^\gamma\sqrt{\frac{\gamma-\nu}{\nu}}
\]
is meaningful only for $0<\nu<\gamma$.  The endpoint $\nu=0$ should be
handled separately.

For $0<\nu<\gamma$, the printed argument is valid.  Indeed, putting
$\theta:=\nu/\gamma$, the displayed choice of $t$ is exactly the value
for which
\[
        s^{-2\nu}
        =
        N_{\theta}^{2}t^{-2\theta}
        \frac{t^2}{t^2+s^{2\gamma}} .
\]
For $\lambda\leq s$ this gives
\[
        s^{-2\nu}
        \leq
        N_{\theta}^{2}t^{-2\theta}
        \frac{t^2}{t^2+\lambda^{2\gamma}},
\]
and hence, using equation (24),
\[
        s^{-2\nu}\norm{E_{[0,s)}x}{}^2
        \leq
        N_{\theta}^{2}t^{-2\theta}|K_t^\gamma(x)|^2
        \leq
        N_{\theta}^{2}\norm{x}{\nu:\gamma}^{2} .
\]
Taking the supremum over $s>0$ yields
\[
        \normiii{x}{\nu}
        \leq
        N_{\nu/\gamma}\norm{x}{\nu:\gamma}
\]
for $0<\nu<\gamma$.

It remains only to record the endpoint $\nu=0$.  Since $N_0=1$, the
claim becomes
\[
        \normiii{x}{0}\leq\norm{x}{0:\gamma}.
\]
In fact equality holds.  First,
\[
        \normiii{x}{0}
        =
        \sup_{s>0}\norm{E_{[0,s)}x}{}
        =
        \norm{x}{}.
\]
The last equality follows because the spectral measure of $T^*T$ is
compactly supported and $E_{[0,s)}$ equals the identity for all sufficiently
large $s$.  On the other hand, by equation (24),
\[
        |K_t^\gamma(x)|^2
        =
        \int_{[0,\infty)}
        \frac{t^2}{t^2+\lambda^{2\gamma}}\,d\mu_x(\lambda)
        \leq
        \int_{[0,\infty)}d\mu_x(\lambda)
        =
        \norm{x}{}^2 .
\]
Thus $\norm{x}{0:\gamma}\leq\norm{x}{}$.  Conversely, as
$t\to\infty$ the integrand in equation (24) increases pointwise to $1$;
by monotone convergence,
\[
        \lim_{t\to\infty}|K_t^\gamma(x)|^2
        =
        \norm{x}{}^2 .
\]
Therefore
\[
        \norm{x}{0:\gamma}
        =
        \sup_{t>0}|K_t^\gamma(x)|
        =
        \norm{x}{}
        =
        \normiii{x}{0}.
\]
This proves inequality (26b) also at $\nu=0$.  Thus the statement of
Proposition 2.2 is unchanged; only the endpoint argument needs
to be supplied separately.


\subsection*{Erratum E6: the endpoint $\nu=0$ in Proposition 2.2}
\label{s:erratum:E6}

In the proof of Proposition 2.2, specifically in the proof of
inequality (26b), the displayed choice
\[
        t=s^\gamma\sqrt{\frac{\gamma-\nu}{\nu}}
\]
is meaningful only for $0<\nu<\gamma$.  The endpoint $\nu=0$ should be
handled separately.

For $0<\nu<\gamma$, the printed argument is valid.  Indeed, putting
$\theta:=\nu/\gamma$, the displayed choice of $t$ is exactly the value
for which
\[
        s^{-2\nu}
        =
        N_{\theta}^{2}t^{-2\theta}
        \frac{t^2}{t^2+s^{2\gamma}} .
\]
For $\lambda\leq s$ this gives
\[
        s^{-2\nu}
        \leq
        N_{\theta}^{2}t^{-2\theta}
        \frac{t^2}{t^2+\lambda^{2\gamma}},
\]
and hence, using equation (24),
\[
        s^{-2\nu}\norm{E_{[0,s)}x}{}^2
        \leq
        N_{\theta}^{2}t^{-2\theta}|K_t^\gamma(x)|^2
        \leq
        N_{\theta}^{2}\norm{x}{\nu:\gamma}^{2} .
\]
Taking the supremum over $s>0$ yields
\[
        \normiii{x}{\nu}
        \leq
        N_{\nu/\gamma}\norm{x}{\nu:\gamma}
\]
for $0<\nu<\gamma$.

It remains only to record the endpoint $\nu=0$.  Since $N_0=1$, the
claim becomes
\[
        \normiii{x}{0}\leq\norm{x}{0:\gamma}.
\]
In fact equality holds.  First,
\[
        \normiii{x}{0}
        =
        \sup_{s>0}\norm{E_{[0,s)}x}{}
        =
        \norm{x}{}.
\]
The last equality follows because the spectral measure of $T^*T$ is
compactly supported and $E_{[0,s)}$ equals the identity for all sufficiently
large $s$.  On the other hand, by equation (24),
\[
        |K_t^\gamma(x)|^2
        =
        \int_{[0,\infty)}
        \frac{t^2}{t^2+\lambda^{2\gamma}}\,d\mu_x(\lambda)
        \leq
        \int_{[0,\infty)}d\mu_x(\lambda)
        =
        \norm{x}{}^2 .
\]
Thus $\norm{x}{0:\gamma}\leq\norm{x}{}$.  Conversely, as
$t\to\infty$ the integrand in equation (24) increases pointwise to $1$;
by monotone convergence,
\[
        \lim_{t\to\infty}|K_t^\gamma(x)|^2
        =
        \norm{x}{}^2 .
\]
Therefore
\[
        \norm{x}{0:\gamma}
        =
        \sup_{t>0}|K_t^\gamma(x)|
        =
        \norm{x}{}
        =
        \normiii{x}{0}.
\]
This proves inequality (26b) also at $\nu=0$.  Thus the statement of
Proposition 2.2 is unchanged; only the endpoint argument needs
to be supplied separately.
